\section{\label{ref:BluetoothSettings}Bluetooth}

Rockbox can send audio to Bluetooth headphones and speakers. This only works
on hardware revision 4 players (Aigo Eros Q V2.2, Hifiwalker H2 V1.9 to V2.3,
Surfans F20 V3.4). On other revisions Bluetooth will refuse to start.

The \setting{Bluetooth Settings} menu is in the main \setting{Settings} menu.

\subsection{Connecting headphones}
Put your headphones in pairing mode, then open \setting{Devices}. Bluetooth is
turned on if it isn't already, and the \dap{} looks for nearby audio devices
for about six seconds. Headphones you have paired before are listed even when
they are switched off.

\begin{description}
  \item[Play.] Connect to the selected device. The first connection also pairs
    with it. The list title shows how far along the connection is.
  \item[Long press Menu.] Disconnect the selected device.
  \item[Menu.] Turn Bluetooth off and leave the list.
  \item[Back.] Leave the list. Bluetooth stays on and stays connected.
\end{description}

The first row, \setting{Scan again}, starts a new search. Each device shows a
second line: the signal strength for devices found in the search,
``Paired'' for devices you have used before, or ``Connected'' with the codec
in use. If your headphones report their battery level it is shown there too.

Once paired, most headphones will reconnect on their own when you take them
out of their case, as long as Bluetooth is on.

\subsection{What your headphones can do}
\begin{itemize}
  \item Their play, pause, next and previous buttons control Rockbox.
  \item Headphones that pause when you take them off (in-ear detection) will
    pause and resume Rockbox.
  \item Playback pauses when the headphones disconnect or go out of range.
  \item If the headphones support it, the volume keys on the \dap{} change the
    volume of the headphones themselves, and the headphones' own volume
    buttons move the Rockbox volume. Otherwise Rockbox lowers the volume
    digitally before sending the audio.
  \item The headphones' battery level is read when they report it. Most
    earbuds do.
  \item Car stereos and other devices with a display get the title, artist,
    album and playing position.
\end{itemize}

Audio is always sent at 44.1\,kHz. Music at other sample rates is resampled
by Rockbox.

\subsection{Settings}
\begin{description}
  \item[Bluetooth.] Turns Bluetooth on or off. Turning it on reconnects the
    last headphones you used. This setting can be put on the Quick Screen or in
    the Shortcuts menu.
  \item[Devices.] Shows the device list described above.
  \item[Enable on Startup.] Turns Bluetooth on when Rockbox starts and
    reconnects the last headphones. This happens in the background, so it
    doesn't slow down booting. Off by default.
  \item[Auto Off When Idle.] Turns Bluetooth off after this many minutes
    without a connected device, to save battery. Set to \setting{Off} if you
    want headphones to be able to reconnect at any time. Default is 10 minutes.
  \item[Preferred Codec.] The codec to use when the headphones support more
    than one. \setting{Auto} picks the best one both sides support, in the
    order LDAC, aptX HD, aptX, SBC. If your choice isn't supported by the
    headphones, the next best is used. SBC is supported by all headphones and
    is the most tested.
  \item[LDAC Quality.] \setting{Adaptive} lowers the bitrate when the
    connection gets weak. The fixed settings are 990, 660 and 330\,kbps.
  \item[Show Bluetooth Icon.] Shows the Bluetooth status and headphone battery
    in the status bar and in themes that support it.
  \item[Forget Devices.] Lists paired devices. Select one and confirm to
    remove it. Bluetooth is turned on if needed, since the list of paired
    devices is kept by the Bluetooth stack.
  \item[Debug Logging.] Writes \fname{/.rockbox/bt\_log.txt} and
    \fname{/.rockbox/bt\_hci.pcap}. Leave it off unless you are tracking down
    a problem, as it writes to the disk often.
\end{description}

\subsection{Device name}
Headphones see the \dap{} as ``Rockbox''. To use another name, put it on the
first line of \fname{/.rockbox/playername.txt}.

\subsection{Files}
\begin{description}
  \item[\fname{/.rockbox/bt\_config.conf}] Paired devices and their keys.
    Don't share this file.
  \item[\fname{/.rockbox/bt\_last.dat}] The last connected device, used for
    reconnecting.
  \item[\fname{/.rockbox/bt\_log.txt}] Event log, only with
    \setting{Debug Logging}.
  \item[\fname{/.rockbox/bt\_hci.pcap}] Capture of the traffic between Rockbox
    and the Bluetooth chip during connection setup and the first seconds of
    playback, only with \setting{Debug Logging}. It can be opened with
    Wireshark.
\end{description}

\subsection{Troubleshooting}
\begin{description}
  \item[``Bluetooth failed to start''.] The message names the step that
    failed. ``controller'' means the Bluetooth chip did not answer. Try again
    after a restart.
  \item[Headphones are not in the list.] Make sure they are in pairing mode,
    then use \setting{Scan again}. Only audio devices are listed.
  \item[``Could not connect''.] The headphones may be connected to another
    device, such as a phone. Disconnect them there first.
  \item[Audio drops out.] Bluetooth needs a clear path between the \dap{} and
    the headphones. Keeping the \dap{} in a pocket on the same side as the
    headphones helps. With LDAC, try \setting{Adaptive} or a lower quality.
\end{description}
